% !TEX root = ../main.tex
\section{What We Found}

\begin{frame}{Seven of Nine Systems Violate Properties}
  \begin{center}
    \footnotesize
    \renewcommand{\arraystretch}{0.95}
    \begin{tabular}{@{}llc@{}}
      \toprule
      \textbf{System} & \textbf{Property violated} & \textbf{Network} \\
      \midrule
      SwitchML  & computational correctness & out-of-order \\
      ATP       & memory leak freedom       & out-of-order \\
      NetReduce & terminality               & lossy \\
      \addlinespace[0.45em]
      NetCache  & cache consistency & lossy \\
      NetCache  & terminality       & reliable \\
      FarReach  & cache consistency & reliable \\
      \addlinespace[0.45em]
      NetLock   & lock exclusion & delayed \\
      FissLock  & lock exclusion & delayed \\
      \bottomrule
    \end{tabular}
  \end{center}
  \vspace{0.7em}
  \begin{center}
    \begin{minipage}{0.94\textwidth}
      \centering
      \visible<2->{Few of these systems explicitly state their network assumptions.}
    \end{minipage}
  \end{center}
  \speechnote{13}
\end{frame}

\begin{frame}{Not Just on Paper}
  \vspace{0.5em}
  \begin{itemize}
    \item Following each violation trace, we injected exactly those losses
          and delays into the \textbf{real open-source implementations}.
    \vspace{0.7em}
    \item<2-> \textbf{All six open-source systems reproduced the violation.}
    \vspace{0.7em}
    \item<3-> The authors of \textbf{ATP} and \textbf{NetReduce} confirmed these
          were oversights at design time.
    \vspace{0.7em}
    \item<4-> \textbf{EPIC} (SIGCOMM~'26) used VeriNC \emph{during} its design ---
          a buffer management pitfall was fixed before publication.
  \end{itemize}
  \speechnote{14}
\end{frame}

\begin{frame}{Three Patterns Behind Those Violations}
  \vspace{0.6em}
  \begin{itemize}
    \item \textbf{Stale packets pollute switch state.}\\
          {Give the switch a way to recognize them.}
    \vspace{0.8em}
    \item \textbf{Concurrent events race with no coordination.}\\
          {Handle how events interact, not one event at a time.}
    \vspace{0.8em}
    \item \textbf{Unreliability turns a latent risk into a real violation.}\\
          {Most of what we found needs only a single such event.}
  \end{itemize}
  \speechnote{15}
\end{frame}

\begin{frame}{Conclusion}
  \vspace{0.6em}
  % 条目间距用 itemsep 统一给，不要逐条插 \vspace——手插的间距在折行的条目
  % 之间表现不一致，三条看着一条比一条挤。preamble 的钩子把 itemsep 设成 2pt，
  % 在环境内部再设一次就能盖掉（钩子先执行）。
  \begin{itemize}
    \setlength{\itemsep}{0.9em}
    \item \textbf{Design-level verification} for in-network computing systems
    \item A high-level \textbf{specification language} and
          \textbf{configurable network environments}
    \item Nine systems modeled; design risks in seven of them, found in
          seconds and reproduced on the real implementations
  \end{itemize}
  \vspace{1.0em}
  \begin{center}
    \begin{minipage}{0.94\textwidth}
      \centering
      \alert{\textbf{Formal modeling and verification are an essential step}}\par
      \alert{\textbf{in driving the deployment of INC systems.}}
    \end{minipage}
  \end{center}
  \speechnote{16}
\end{frame}
